The dielectric response in a semiconducting or insulating material arises when the electric charges in the material cannot be effectively drained away when an external electric field $\bm{E}$ is applied across the material. Instead, the electronic and ionic charges in the material are redistributed, inducing an internal polarisation to counter--balanced the external field. The strength of the (frequency--dependent) dielectric response $\bm{D}(\omega)$ is correlated with the electric field  as:
\begin{equation}\label{dielectric_reponse}
    \bm{D}(\omega)=\varepsilon(\omega)\bm{E}(\omega),
\end{equation}
where  $\varepsilon(\omega)$ is the tensor of dielectric permittivity that is an intrinsic property of the material. Generally for a complex material, its dielectric response is contributed by multiple different mechanisms, including (in the order of decreasing magnitudes) space--charge ($\varepsilon_{sc}$), dipole reorientation ($\varepsilon_{dip}$), electronic (or optical, $\varepsilon_{elec}$) and ionic (or phonon, $\varepsilon_{ion}$) contributions. Based on the \emph{modern theory of polarisation},\cite{spaldin2012beginner} algorithms for calculating the latter two quantities in condensed matters have been implemented in many modern electronic structure codes. Therefore, the discussions in this section will be focusing on the electronic and ionic contributions to the polarisations.

\paragraph{\textbf{Electronic polarisation}}

In molecular systems, the dipole moment for an solated molecule is trivial to define and understand.\cite{stone2013theory} For example, a simple linear molecule of carbon monoxide (CO) is said to carry a \textit{permanent dipole} moment because of the difference in the electronegativity between the C and O atom, leading to a partial point charge of $+\delta$ and $-\delta$ to be possessed by each atom, respectively. The corresponding dipole moment for CO is therefore $\bm{\mu}=\delta\mathbf{r}/|\mathbf{r}|$, where $\mathbf{r}$ is the vector joining the two nuclei in the molecule. Within this point charge model, molecules such as O$_2$ is said to be nonpolar for its homoatomic nature. If we further take  the quantum mechanical nature of the electrons into account, then we need to think of the electrons in the molecules as a spatial distribution of charge densities around the nuclei, regardless of whether the molecule is polar or nonpolar. Firstly, when an electric field is applied to the molecule, an \textit{induced dipole} on the molecule is generated resulted from the distortion in the shape of the charge density as a response to the external field. Secondly, a molecule always possess an \textit{instantaneous dipole} due to the non--static nature of the electron. The constant spatial movements of electrons causes a fluctuation in the spatial distribution of the electronic charge density with respect to the positions of the nuclei, meaning at any time instance $t$, there will be some finite net charge of $\delta_{t,i}$ associated with the $i$--th atom in the molecule, thus gives rise to a dipole moment in the molecule. Further considering the electronic correlation effects, the fluctuations of the electron densities between two nearby molecules\cite{dobson2001prediction} or nanostructures\cite{tawfik2018evaluation} are correlated, which contributes to the important London dispersion interactions\cite{stone2013theory} between two non--covalently bounded entities.

However, such a simple exercise in counting the ionic charges and interatomic distances to determine the dipole polarisations fails for extended solids,\citep{spaldin2012beginner} which leads to the electronic polarisation become multi--valued that is dependent on the choice of the lattice vector. This problem can be overcome by making a connection to the experiments, in which the \emph{change in the polarisation} $\Delta\mathbf{P}_{elec}$ between the initial  ($\lambda_i$) and final ($\lambda_f$) state in a sample, defined as\cite{resta1992theory}
\begin{equation}\label{polarisation_change}
\Delta\mathbf{P}_{elec}=\int_{\lambda_i}^{\lambda_f}\frac{\partial \mathbf{P}_{elec}}{\partial \lambda}d\lambda,
\end{equation}
is measured as an electric current flowing through a closed circuit connecting the two ends of the sample. Here, $\lambda$ is the coordinate in the phase space that can be spanned by time, ionic coordinates, external electric and/or magnetic fields, \emph{etc}, that drive the change in $\mathbf{P}_{elec}$ along a path that goes from $\lambda_i$ to $\lambda_f$ in the phase space. For simplicity here, we consider $\lambda$ is represented by the time variable, such that $\frac{\partial \mathbf{P}_{elec}}{\partial \lambda}=\frac{\partial \mathbf{P}_{elec}}{\partial t}$. This means now, in order to obtain the change of polarisation over time, we need to take into account the time--dependent evolution of a wavefunction $\psi(\bm{r}(t))$ under a time--dependent Hamiltonian $\mathcal{H}(\bm{r}(t))$, in which $\partial \bm{P}/\partial t$ can be correlated with the change in the electronic wavefunction over $\Delta t$, and the latter can be obtained from the perturbation theory as an expansion involving the operator $\partial \mathcal{H}(\bm{r}(t))/\partial t$. However, note that in the ground state electronic structure theory, \emph{e.g.} DFT, the time--dependency is not included. This problem is overcome under the theory of Berry phase with the recognition of the chain rule for differentiation $\partial \mathcal{H}(\bm{r}(t))/\partial t=\bm{\nabla}_{\bm{r}}\mathcal{H}\cdot\dot{\bm{r}}$, such that \cref{polarisation_change} can be cast into a form that is  only dependent on the atomic coordinates $\bm{r}$ but not time. For periodic DFT, we further take into account the KS orbitals in the form of the Bloch wave $\psi_{n,k}(\bm{r})=e^{i\bm{k}\bm{r}}u_{n,k}(\bm{r})$, we arrive at the  following workable definition of electronic polarisation that depends only on the occupied electronic states:
\begin{equation}\label{electron_pol_DFT}
\bm{P}_{elec}=\frac{e}{(2\pi)^{3}}\int_{\bm{k}\in BZ}\mathcal{A}(\bm{k})\mbox{d}\bm{k},
\end{equation}
in which $\mathcal{A}(\bm{k})=-i\langle u_{n,k}|\bm{\nabla}_{\bm{k}}|u_{n,k}\rangle$ is the so--called Berry curvature.  Based on \cref{dielectric_reponse}, the corresponding electronic dielectric tensor can be obtained from:\cite{souza2002first}
\begin{equation}\label{eq:electronic_dielectric_tensor}
    \varepsilon^{\alpha\beta}_{elec}(\omega)=\delta_{\alpha\beta}+4\pi\frac{\partial \bm{P}_{\alpha}}{\partial \bm{E}_{\beta}},
\end{equation}
implying that to determine $\varepsilon_{elec}(\omega)$ algorithmically, the polarisation tensor needs to be evaluated under the application of different electric field, from which the partial derivative in \cref{eq:electronic_dielectric_tensor} can be evaluated. 

\paragraph{\textbf{Ionic polarisation}} The above discussions has highlighted the approach to calculate the electronic polarisation in a material with fixed ionic positions $\bm{R}_{0}$, such as the equilibrium structure that corresponds a local minimum on the PES. On the other hand, the ionic polarisation $\bm{P}_{ion}$ reflects the change in the system polarisation resulted from the variation in the  charge density distribution when the ions change their positions from $\bm{R}_{0}$ to $\bm{R}$ along a particular normal mode $\bm{\mathcal{Q}}$ of lattice vibrations. This is the case because under the Born--Oppenheimer approximation,  the electron density is assumed to change its shape instantaneously with respect to the change in the ionic positions. It also implies that evaluating $\bm{P}_{ion}$ requires the knowledge of the vibrational eigenmodes and eigenfrequencies $\{\bm{\mathcal{Q}}_n,\omega_n\}$, together with  $\bm{P}_{elec}$ to be determined for different $\bm{R}$. 

To map the above interpretation back to a more classical picture of dipole moments (\textit{c.f.} the discussions about molecular dipoles), the Born effective charge tensor on ion $i$ in a solid is defined\cite{ghosez1998dynamical} as the change in the polarisation along the $\alpha$--Cartesian direction when the ion is displaced along the normal mode in the $\beta$--direction at zero field:
\begin{equation}\label{eq:born_charge}
    Z_{i,\alpha\beta}^{*}=\left.\frac{\Omega}{|e|}\frac{\partial \bm{P}_{\alpha}}{\partial \bm{\mathcal{Q}}_\beta^{(i)}}\right\vert_{E=0},
\end{equation}
where $\Omega$ is the unit cell volume. With this, the ionic contribution to the polarisation in a crystal unit cell can be calculated as:
\begin{equation}\label{eq:ionic_pol}
    \Delta\bm{P}_{ion}=\frac{|e|}{\Omega}\sum_{i}\bm{Z}_{i}^{*}\bm{\mathcal{Q}}^{(i)},
\end{equation}
where the summation is performed all ions in the unit cell. The ionic contribution to the tensor of dielectric permittivity is then given by:\cite{gonze1997dynamical}
\begin{equation}\label{eq:ionic_epsilon}
     \varepsilon_{ion}^{\alpha\beta}(\omega)=\frac{4\pi}{\Omega}\sum_{n=1}^N\frac{\left[\sum_{s\alpha^\prime}Z_{s,\alpha\alpha^\prime}^{*}\bm{\mathcal{Q}_n^{*}}\right] \left[\sum_{s^\prime\beta^\prime}Z_{s^\prime,\beta\beta^\prime}^{*}\bm{\mathcal{Q}}_n\right]}{\omega_n^2-\omega^2},
\end{equation}
and the index $n$ runs over all phonon eigenstates. The concept of the Born effective charge has been a particularly useful property indicator to explain and screen for the emergence of spontaneous polarisation in materials, such as $d^0$ transition metal perovskites,\cite{matthias1949new} in which the $d^0$ B--site cation exhibits an anomalously large Born effective charge compared to its formal charge.\cite{armiento2011screening} 


